\section{Applications}\label{sec:applications}

\subsection{Gaming}\label{sec:applications_gaming}
Multiplayer video games extensively use random mechanics that directly impact player experiences and in-game economies. From critical hit calculations to loot drops and lootbox mechanics, players invest significant time and resources based on these probabilities. The fairness of these systems directly affects player satisfaction and game balance.

This need for verifiable randomness extends naturally into gambling applications, where every slot machine pull, card shuffle, and dice roll must be demonstrably random and tamper-proof to maintain player trust and regulatory compliance. The stakes are often significant, making manipulation attempts a serious concern.

\subsection{NFTs}\label{sec:applications_nfts}
Non-Fungible Tokens (NFTs), such as those from the Bored Ape Yacht Club (BAYC) \cite{boredapeyachtclub}, are typically generated by combining various image layers to create unique digital collectibles. Due to the challenges associated with generating randomness directly on-chain—a situation often referred to as the "On-Chain Randomness Problem" —these NFTs are frequently created off-chain by developers and then uploaded to the blockchain. This process can be susceptible to manipulation. Implementing true on-chain randomness would enable provably fair and transparent NFT minting processes, ensuring each token's uniqueness and integrity.
\subsection{Proof of Stake Consensus}\label{sec:applications_defi}
Randomness is extensively utilized in Proof of Stake (PoS) consensus mechanisms.  For instance, Solana employs random selection to determine the next leader in its consensus process \cite{yakovenko2025solana}. This method ensures a fair and unpredictable rotation among validators.
\subsection{Artificial Intelligence}\label{sec:applications_security}
Random values are essential in neural networks such as large language models (LLMs) like ChatGPT because they help the system learn. One of the most important uses of randomness is weight initialization, where the initial values of connections (weights) are set unpredictably. If weights started from known values, an attacker or even the system itself could predict learning paths, reducing diversity and making the network less effective. Instead, true random initialization ensures that different neurons learn different features, making the network more powerful and capable of recognizing complex patterns. In machine learning, unpredictability is key to avoiding bias and ensuring robust, secure models \cite{pmlr-v9-glorot10a}.
